% 1405/6/13
% H. Behrouz 
% Email : Texlivegroup@gmail.com


\documentclass[a4paper,14pt]{article}
\usepackage[a4paper,vmargin={20mm,20mm},hmargin={10mm,20mm}]{geometry}
\usepackage{fancybox}
\usepackage{array}
\usepackage{amsmath,amssymb}
\usepackage{mathtools}


\usepackage[colorlinks,linkcolor=blue,citecolor=blue]{hyperref}
\usepackage{wasysym}
\usepackage{amsmath} 
\usepackage{amssymb}
\usepackage{amsfonts}
\usepackage{mathtools}
\usepackage{float}


\usepackage{xepersian}
\settextfont[Scale=1.3]{Yas}
\setdigitfont[Scale=1]{Yas}
\setlatintextfont[Scale=1.25]{Times New Roman}




\setlength{\fboxsep}{0.5cm}
\renewcommand{\baselinestretch}{2}
\DeclareMathSizes{10}{13.5}{7}{7}








\begin{document}
%دستوری برای ایجاد کادر صفحه
\fancypage{\setlength{\fboxrule}{0.5mm}\fbox}{}	
	
\thispagestyle{empty}	
\begin{center}
\textbf{
 بنام خدا
}
\end{center}
\vspace*{-1.5cm}
\begin{flushleft}
تاریخ: 
$1405/8/20$
\end{flushleft}
\textbf{
نام مدرسه:
}	
\hspace{11cm}
\textbf{
نام درس:
}
\\
\textbf{
نام و نام خانوادگی:
}	
\hspace{3.5cm}
\textbf{
نام آموزگار:
}
\hspace{3.5cm}
\textbf{
سال تحصیلی:
}
\\		 
\noindent\makebox[\textwidth]{\rule{19cm}{1pt}}
\\	
$1$-	
هرگاه داشته باشیم 
$x^{2}+y^{2}=kxy$
و 
$k>2$
و 
$x>y>0$ 
باشد 
$\dfrac{x+y}{x-y}$
برابر چیست؟ 
\vspace*{-1.3cm}
\begin{flushleft}
\textbf{بارم}
\end{flushleft}
\noindent\makebox[\textwidth]{\rule{19cm}{1pt}}
\\		
$2$-
ثابت کنید 
$(ax+by)^{2}+(bx-ay)^{2}=(a^{2}+b^{2})(x^{2}+y^{2})$	
و از آنجا نتیجه بگیرید چنانچه 
$ax+by=k$
باشد 
$x^{2}+y^{2}$
وقتی مینیمم است که 
$\dfrac{x}{a}=\dfrac{y}{b}$
باشد.
\vspace*{-1.3cm}
\begin{flushleft}
\textbf{بارم}
\end{flushleft}	 
\noindent\makebox[\textwidth]{\rule{19cm}{1pt}}
\\	
$3$-
اگر 
$\dfrac{1}{x}+\dfrac{1}{y}+\dfrac{1}{z}=1$
و 
$ax^{3}=by^{3}=cz^{3}$
باشد ثابت کنید:
\vspace*{-1.3cm}
\begin{flushleft}
\textbf{بارم}
\end{flushleft}
\begin{flushleft}
$\sqrt[3]{ax^{2}+by^{2}+cz^{2}}=\sqrt[3]{a}+\sqrt[3]{b}+\sqrt[3]{c}$
\end{flushleft}		
\noindent\makebox[\textwidth]{\rule{19cm}{1pt}}
\\
$4$-
عبارت‌های 
$a^{5n}+a^{n}+1$
و 
$a^{4n}+a^{2n}+1$
را تجزیه کنید و بزرگ‌ترین مقسوم علیه مشترک آنها را به 
دست آورید.
\\
\vspace*{-1.3cm}
\begin{flushleft}
\textbf{بارم}
\end{flushleft}
\noindent\makebox[\textwidth]{\rule{19cm}{1pt}}
\\
$5$-
ثابت کنید از برابری 
$\dfrac{x}{a}=\dfrac{y}{b}=\dfrac{z}{c}$
نتیجه می‌شود:
\vspace*{-1.3cm}
\begin{flushleft}
\textbf{بارم}
\end{flushleft}
\begin{flushleft}
$\left(a^{2}+b^{2}+c^{2} \right)\left(x^{2}+y^{2}+z^{2} \right)=(ax+by+cz)^{2}  $
\end{flushleft}
\noindent\makebox[\textwidth]{\rule{19cm}{1pt}}
\\
$6$-
حاصل عبارت 
$A=\left(1+\dfrac{1}{2} \right)+\left(1+\dfrac{1}{4} \right)+\left(1+\dfrac{1}{16} \right)+\cdots+\left(1+\dfrac{1}{2^{2^{n}}} \right)+ $ 
چیست؟
\vspace*{-1.3cm}
\begin{flushleft}
\textbf{بارم}
\end{flushleft}
\noindent\makebox[\textwidth]{\rule{19cm}{1pt}}
\\
$7$-
هرگاه داشته باشیم 
$\sqrt[3]{x^{2}}+\sqrt[3]{x}=1$
حاصل 
$\dfrac{1-x-x^{2}}{3x}$
برابر چیست؟
\vspace*{-1.3cm}
\begin{flushleft}
\textbf{بارم}
\end{flushleft}
\noindent\makebox[\textwidth]{\rule{19cm}{1pt}}
\\
$8$-
ثابت کنید:
\vspace*{-1.3cm}
\begin{flushleft}
\textbf{بارم}
\end{flushleft}
\begin{center}
$\sqrt{2}\left(\sqrt[4]{7+4\sqrt{3}}-\sqrt[4]{7-4\sqrt{3}} \right)=2 $
\end{center}
\noindent\makebox[\textwidth]{\rule{19cm}{1pt}}
\\
$9$-
در معادله 
$x^{2}-2x\sin\alpha+1=0$
حاصل 
$\dfrac{1}{x^{\prime}4}+\dfrac{1}{x^{\prime\prime}4}$
برابر چیست؟
\vspace*{-1.3cm}
\begin{flushleft}
\textbf{بارم}
\end{flushleft}
\noindent\makebox[\textwidth]{\rule{19cm}{1pt}}
\\
\begin{center}
«موفق باشید»
\end{center}
\end{document}